\Needspace{20\baselineskip}
\section{Estado del Arte}

\subsection{Avances y limitaciones recientes en Text-to-SQL}

La literatura reciente destaca una evolución significativa en la tarea de
Text-to-SQL, impulsada por nuevos marcos de trabajo que mejoran la alineación
entre la intención del usuario y el esquema de la base de datos
\citep{qin2024gensql}. Estos avances se apoyan en la capacidad de los grandes
modelos de lenguaje para procesar semántica compleja, superando ampliamente a
los enfoques tradicionales basados en reglas heurísticas o análisis
sintáctico superficial \citep{luo2024nl2sql}.

En términos metodológicos, arquitecturas recientes como Gen-SQL proponen
mecanismos eficientes para acortar la brecha entre el lenguaje natural y la
base de datos, optimizando la identificación de tablas y columnas relevantes
previo a la generación del código \citep{qin2024gensql}. De manera
complementaria, marcos de trabajo industriales integran procesos estandarizados
que buscan garantizar la consistencia sintáctica y semántica de las consultas
generadas, un factor crítico para su adopción en ecosistemas de datos
productivos \citep{gladkykh2025datrics}.

A pesar de este progreso empírico, la resolución de consultas analíticas que
involucran múltiples cruces (\textit{joins}), subconsultas anidadas y
operaciones matemáticas sigue siendo un problema abierto \citep{luo2024nl2sql}.
La introducción de entornos de evaluación rigurosos, como el \textit{benchmark}
Spider 2.0, expone estas vulnerabilidades al someter a los modelos a tareas
realistas y altamente desafiantes. Los resultados de estas pruebas demuestran
que la precisión de ejecución (\textit{Execution Accuracy}) cae drásticamente
en esquemas masivos, lo que subraya la necesidad de integrar mecanismos de
validación externa, como la orquestación de agentes, para asegurar la fiabilidad
operativa \citep{lei2024spider}.

\subsection{Despliegue de RAG y agentes en entornos corporativos}

El paso de entornos controlados de laboratorio a implementaciones productivas
representa uno de los focos principales de la investigación actual en
Inteligencia de Negocios. Para que los modelos generativos sean viables
corporativamente, la literatura subraya la necesidad de acoplar técnicas
de recuperación de contexto que limiten el conocimiento a dominios
específicos y mitiguen las alucinaciones \citep{bourdin2024modelos}.

Estudios empíricos recientes documentan la efectividad de la Generación
Aumentada por Recuperación (RAG) en escenarios de alta criticidad. Por ejemplo,
la implementación de estas arquitecturas en instituciones financieras ha
demostrado mejoras sustanciales en la gestión eficiente de \textit{insights},
permitiendo consultas seguras sobre repositorios documentales y esquemas
de bases de datos sin exponer la privacidad de la información
\citep{arguello2025rag}.

Sin embargo, el despliegue de sistemas orquestados por agentes de IA enfrenta
obstáculos operativos significativos, especialmente en PyMEs y ecosistemas
de la Industria 4.0 \citep{bourdin2024modelos}. Las investigaciones
advierten que las barreras de infraestructura, la falta de gobernanza de datos
y la complejidad de gestionar múltiples agentes autónomos dificultan la
estabilidad del sistema \citep{maddila2024agentes}. Esto evidencia que el
desafío actual no radica únicamente en la generación de la respuesta, sino
en la arquitectura de despliegue que sostiene al agente interactuando
con el entorno empresarial.

\subsection{Fronteras en la visualización automatizada de datos}

En el dominio de la representación gráfica, los enfoques más vanguardistas
aprovechan la capacidad de los grandes modelos de lenguaje para ir más allá de
la consulta tabular y traducir comandos en lenguaje natural directamente hacia
visualizaciones. El objetivo principal de estos esfuerzos es automatizar el
ciclo final del Business Intelligence, permitiendo a usuarios no técnicos
consumir información en formatos gráficos de manera interactiva
\citep{perez2026ia}.

Para lograr una traducción efectiva entre el lenguaje natural y las gramáticas
de visualización, la investigación actual emplea metodologías de razonamiento
escalonado (\textit{Chain-of-Thought}). Frameworks como DeepVIS demuestran
que obligar al modelo a planificar lógicamente el tipo de gráfico, los ejes
y las agregaciones antes de generar el código reduce drásticamente la
brecha semántica y mejora la precisión del diseño \citep{shuai2025deepvis}.
Paralelamente, surge una fuerte tendencia hacia la creación de arquitecturas
ligeras, diseñadas específicamente para generar tableros completos sin la
sobrecarga de herramientas de BI tradicionales
\citep{shi2026nl2dashboardlightweightcontrollableframework}.

No obstante, la literatura advierte que la utilidad real de estos sistemas
aún está en fase de maduración. Las alucinaciones en la escala de los ejes,
la selección incorrecta de paletas de colores y la falta de interactividad
avanzada son problemas recurrentes. Además, la fidelidad gráfica comparativa
sigue siendo un obstáculo crítico, limitando la confianza del usuario final
en las interfaces generadas íntegramente por Inteligencia Artificial
\citep{perez2026ia}.

\subsection{Desafíos abiertos en la evaluación comparativa}

A pesar de los rápidos avances en la generación de código y visualizaciones,
la evaluación de sistemas de Inteligencia de Negocios Generativa (GenBI)
constituye uno de los mayores cuellos de botella metodológicos. Las
aproximaciones clásicas, basadas en métricas deterministas, resultan
insuficientes para capturar los matices del razonamiento analítico, ya que
una respuesta válida en BI puede presentarse con múltiples variaciones
estructurales y visuales \citep{perez2026ia}.

Para abordar este problema, investigaciones publicadas entre 2024 y 2026
enfatizan la necesidad de ejecutar análisis comparativos exhaustivos sobre
el rendimiento de distintos modelos fundacionales, evaluando su capacidad
para operar dentro de flujos de trabajo analíticos reales \citep{Jaisinghani2024}.
En este contexto, el uso de LLMs como jueces automatizados se postula
como la técnica más prometedora para evaluar semánticamente tanto el código
SQL generado como la calidad y coherencia de las representaciones gráficas.

A pesar de estos esfuerzos, la carencia de marcos de evaluación integrales
que auditen el ciclo de vida completo de los datos —desde la recuperación
del esquema y la orquestación multi-agente, hasta el diseño final del
tablero— representa una brecha crítica en el estado del arte. Esta limitación
documentada subraya la relevancia de proponer nuevas arquitecturas
supervisadas que garanticen la trazabilidad y la corrección técnica en
entornos analíticos generativos.